\chapter{Hemorrhoids}
\index{Hemorrhoids}
\label{sec:Hemorrhoids}

\section{Overview}
Hemorrhoids are swollen, enlarged veins in the lower rectum and anal canal that can cause bleeding, pain, and discomfort during bowel movements. They affect approximately 50\% of adults over the age of 50, with peak prevalence between ages 45 and 65.
\section{Classification}

\begin{description}[font=\normalfont\bfseries,leftmargin=3em,labelwidth=2.5em]
  \item[Cause] Vascular/Age-related
  \item[Organ System] Gastrointestinal
  \item[Pathophysiology] Increased venous pressure in hemorrhoidal plexuses leading to vascular dilation, swelling, and prolapse of anal cushions
  \item[Duration] Acute or Chronic
  \item[Onset] Gradual
  \item[Mode of Transmission] Non-communicable
  \item[Clinical Course] Chronic, relapsing
  \item[Severity] Mild to Moderate
  \item[Extent of Disease] Localized
  \item[Age Group] Adults
  \item[Epidemiology] Affects up to 50\% of adults; most common in ages 45--65
  \item[ICD-11 Code] DB60.0
\end{description}

\section{Causes}
Hemorrhoids develop when increased pressure in the lower rectum and anal canal causes the vascular cushions to swell and distend, most commonly from chronic straining during bowel movements. Prolonged sitting on the toilet, chronic constipation or diarrhea, and low-fiber diets contribute to venous engorgement of the internal and external hemorrhoidal plexuses. This condition develops from a combination of genetic and environmental factors that disrupt normal body processes. When these normal processes are disrupted, it leads to the symptoms and signs of the condition. Several factors can work together to trigger or make the condition worse. Knowing the specific cause helps your doctor choose the most effective treatment.

\section{Risk Factors}

\begin{itemize}[noitemsep]
  \item Chronic constipation and straining during defecation
  \item Low-fiber diet
  \item Prolonged sitting on the toilet
  \item Pregnancy and vaginal childbirth
  \item Obesity
\end{itemize}

\section{Signs and Symptoms}

\subsection{Early symptoms}
\begin{itemize}[noitemsep]
  \item Early manifestations of hemorrhoids may be subtle and nonspecific
\end{itemize}

\subsection{Common symptoms}
\begin{itemize}[noitemsep]
  \item \subsection{Common symptoms}
  \item \textbf{Common symptoms:}
  \item Painless bright red rectal bleeding, typically seen on toilet paper or dripping into the toilet bowl after a bowel movement
  \item Anal itching, irritation, and a feeling of incomplete evacuation
  \item Prolapsed tissue protruding from the anus, especially with internal hemorrhoids
  \item Pain and swelling around the anus, particularly with thrombosed external hemorrhoids
  \item Pain or discomfort in the affected area
  \item Difficulty performing daily activities
\end{itemize}

\subsection{Less common symptoms}
\begin{itemize}[noitemsep]
  \item Atypical or less common presentations of hemorrhoids may occur
\end{itemize}

\subsection{Emergency warning signs}
\begin{itemize}[noitemsep]
  \item Severe or worsening symptoms of hemorrhoids require immediate medical evaluation
\end{itemize}
\section{Progression and Stages}
Internal hemorrhoids are graded on a scale of 1 to 4 based on the degree of prolapse. Grade 1 hemorrhoids do not prolapse, while Grade 4 hemorrhoids are irreducible and chronically prolapsed. External hemorrhoids are classified separately as thrombosed or non-thrombosed based on the presence of an acute clot.

\section{Complications}

\begin{itemize}[noitemsep]
  \item Thrombosis of an external hemorrhoid causing severe acute pain
  \item Anemia from chronic blood loss
  \item Strangulation of a prolapsed internal hemorrhoid
  \item Reduced quality of life
  \item Emotional and psychological effects
\end{itemize}

\section{Diagnosis}

\begin{itemize}[noitemsep]
  \item Visual inspection of the perianal area and digital rectal examination
  \item Anoscopy for direct visualization of internal hemorrhoids
  \item Lab tests to check how your organs are working
  \item Imaging tests
\end{itemize}

\section{Prevention}

\begin{itemize}[noitemsep]
  \item Eat a high-fiber diet with plenty of fruits, vegetables, and whole grains
  \item Stay hydrated and avoid straining during bowel movements
  \item Go for regular check-ups so any problems are caught early
  \item Follow your doctor's screening recommendations
\end{itemize}

\section{Treatment Options}

\begin{itemize}[noitemsep]
  \item Conservative management with dietary fiber supplementation (psyllium), increased water intake, warm sitz baths, and topical analgesics
  \item Rubber band ligation for Grade 1--3 internal hemorrhoids
  \item Surgical hemorrhoidectomy for Grade 4 or thrombosed hemorrhoids refractory to office procedures
  \item Lifestyle changes and self-care
  \item Surgery when needed
\end{itemize}

\section{Living with It}

\begin{itemize}[noitemsep]
  \item Follow your treatment plan as prescribed
  \item Keep up with regular appointments with your healthcare provider
  \item Adjust your daily habits as needed
  \item Reach out to family, friends, or support groups for help
  \item Keep learning about your condition and how to manage it
\end{itemize}

\section{Outlook}
Most people with hemorrhoids respond well to conservative measures and lifestyle modifications. Your outlook depends on the specific condition, how advanced it is when found, and how well it responds to treatment. Getting treated early generally leads to better results. Many people manage this condition effectively with the right treatment. Regular check-ups are important to monitor your condition and adjust treatment as needed.
